\subsection{素数寄存器实现：Fibonacci 估值、折叠正规形与外置账本}\label{subsec:emergent-arithmetic-prime-register-fib-valuation}

前两小节已经给出稳定算术的抽象骨架：尺度计数配置构成自由交换幺半群，Fibonacci 关系在其上生成同余，而 Zeckendorf 规范形给出每个同值类的唯一横截面。
本小节把这一抽象结构落实为一个完全显式、且可直接计算的素数寄存器实现。

其基本思想是：以一列固定素数充当尺度寄存器，把每个尺度上的计数写成对应素数的指数。
如此一来，尺度计数单子便获得了一个乘法幺半群实现，而 Fibonacci 值映射则转化为指数向量上的加权估值。
于是，前两小节的同值轨道、规范横截面与归一化重写，都可改写为寄存器上的局部指数变换。

与此同时，必须区分两种本质不同的线性化：
寄存器内部的乘法对应的是值的加法叠加；值的乘法则不能由同一内部操作直接实现，而必须经由对象层值本身的普通素因子分解，被外置到一个独立账本中。
因此，本小节得到的不是把一切运算压进同一坐标系的单层编码，而是一个严格分层的实现框架：Fibonacci 寄存器负责加法侧的折叠与规范化，普通素数账本负责乘法侧的外置线性化。

\begin{definition}[素数寄存器单子、Fibonacci 估值与 Zeckendorf 正规寄存器]\label{def:prime-register-fib-valuation-zg-code}
取一列严格递增的素数
\[
\Pi=(\pi_1,\pi_2,\pi_3,\dots),
\]
并定义寄存器状态集合
\[
\mathcal R_{\Pi}
:=
\left\{
\prod_{k\ge1}\pi_k^{a_k}
\;:\;
a_k\in\NN,\ \text{仅有限个 }a_k\neq 0
\right\}.
\]
在通常整数乘法下，\(\mathcal R_{\Pi}\) 构成一个自由交换幺半群，其单位元为 \(1\)。

定义幺半群同构
\[
\Gamma:\mathcal D\to\mathcal R_{\Pi},
\qquad
\Gamma(a):=\prod_{k\ge1}\pi_k^{a_k},
\]
其中 \(\mathcal D\) 为定义 \ref{def:fib-congruence} 之前引入的尺度计数单子。

在 \(\mathcal R_{\Pi}\) 上定义 Fibonacci 估值
\[
\mathrm{Val}_{\mathrm{pr}}:\mathcal R_{\Pi}\to\NN,
\qquad
\mathrm{Val}_{\mathrm{pr}}\!\left(\prod_{k\ge1}\pi_k^{a_k}\right)
:=
\sum_{k\ge1}a_kF_{k+1}.
\]
等价地，
\[
\mathrm{Val}_{\mathrm{pr}}=\mathrm{Val}\circ\Gamma^{-1}.
\]

对任意 \(N\in\NN\)，记其 Zeckendorf 规范形为
\[
Z(N)=(z_1(N),z_2(N),\dots)\in X_{\infty},
\]
并定义相应的规范寄存器
\[
R_F(N):=\prod_{k\ge1}\pi_k^{z_k(N)}.
\]
同时记
\[
\mathsf g(N):=R_F(N),
\qquad
\mathrm{ZG}:=\mathsf g(\NN)\subset\mathcal R_{\Pi}.
\]
\leanverified{valPr}
\leanverified{valPr\_zero}
\leanverified{valPr\_single}
\leanverified{valPr\_add}
\leanverified{paper\_prime\_register\_fib\_valuation\_def}
\end{definition}

\begin{proposition}[尺度计数配置与寄存器状态的同构]\label{prop:prime-register-monoid-realization}
\leanverified{paper\_prime\_register\_monoid\_realization}
映射 \(\Gamma\) 是交换幺半群同构。
因此，尺度计数配置与寄存器状态只是同一自由交换结构的两种表示：前者是指数向量表示，后者是素数幂积表示。
\end{proposition}

\begin{proof}
由算术基本定理，\(\mathcal R_{\Pi}\) 中每个元素都唯一写成 \(\prod_{k\ge1}\pi_k^{a_k}\) 的形式，并且乘法对应指数逐点相加。
因此 \(\Gamma\) 既单射又满射，且满足
\[
\Gamma(a+b)=\Gamma(a)\Gamma(b).
\]
故其为交换幺半群同构。证毕。
\end{proof}

\begin{proposition}[寄存器内部乘法对应值的加法]\label{prop:prime-register-internal-product-adds-values}
对任意 \(r,s\in\mathcal R_{\Pi}\)，有
\[
\mathrm{Val}_{\mathrm{pr}}(rs)=\mathrm{Val}_{\mathrm{pr}}(r)+\mathrm{Val}_{\mathrm{pr}}(s),
\qquad
\mathrm{Val}_{\mathrm{pr}}(1)=0.
\]
因此，寄存器空间内部的乘法并不实现值的乘法，而实现值的加法叠加。
\end{proposition}

\begin{proof}
设
\[
r=\prod_{k\ge1}\pi_k^{a_k},
\qquad
s=\prod_{k\ge1}\pi_k^{b_k}.
\]
则
\[
rs=\prod_{k\ge1}\pi_k^{a_k+b_k},
\]
从而
\[
\mathrm{Val}_{\mathrm{pr}}(rs)
=\sum_{k\ge1}(a_k+b_k)F_{k+1}
=\sum_{k\ge1}a_kF_{k+1}+\sum_{k\ge1}b_kF_{k+1}
=\mathrm{Val}_{\mathrm{pr}}(r)+\mathrm{Val}_{\mathrm{pr}}(s).
\]
而 \(1\) 的全部指数皆为零，故 \(\mathrm{Val}_{\mathrm{pr}}(1)=0\)。证毕。
\leanverified{paper\_prime\_register\_internal\_product\_adds\_values}
\end{proof}

\begin{remark}[寄存器内部乘法的对象层语义]\label{rem:prime-register-internal-product-semantics}
命题 \ref{prop:prime-register-internal-product-adds-values} 表明，寄存器乘法的对象层语义是把两个尺度计数配置叠加到一起。
因此，若 \(R_F(m)\) 与 \(R_F(n)\) 分别表示对象层值 \(m,n\) 的规范寄存器形，则其普通乘积对应的是 \(m+n\) 的未归一化表示，而不是 \(mn\) 的表示。
这说明 Fibonacci 寄存器天然适配加法侧的结构，而不直接适配乘法侧。
\end{remark}

\begin{definition}[寄存器上的局部值保持变换]\label{def:prime-register-local-value-preserving-moves}
在 \(\mathcal R_{\Pi}\) 上定义以下部分可逆局部变换：
\begin{enumerate}
  \item 基准变换
  \[
  \pi_2 \longleftrightarrow \pi_1^2;
  \]
  \item Fibonacci 递推变换：对每个 \(k\ge1\)，
  \[
  \pi_{k+2} \longleftrightarrow \pi_{k+1}\pi_k.
  \]
\end{enumerate}
更一般地，若某寄存器状态包含左端因子，则可把其替换为右端因子；若包含右端因子，则亦可反向替换。
由这些局部部分双射生成的群胚记为 \(\mathcal G_{\Pi}\)。
\end{definition}

\begin{proposition}[寄存器局部变换保持 Fibonacci 估值]\label{prop:prime-register-local-moves-preserve-valuation}
\(\mathcal G_{\Pi}\) 的每一条箭头都保持 \(\mathrm{Val}_{\mathrm{pr}}\) 不变。
也就是说，若 \(r\to r'\) 是由定义 \ref{def:prime-register-local-value-preserving-moves} 的生成元有限复合得到的箭头，则
\[
\mathrm{Val}_{\mathrm{pr}}(r)=\mathrm{Val}_{\mathrm{pr}}(r').
\]
\end{proposition}

\begin{proof}
对生成元逐个验证即可。
对基准变换，
\[
\mathrm{Val}_{\mathrm{pr}}(\pi_2)=F_3=2F_2=\mathrm{Val}_{\mathrm{pr}}(\pi_1^2).
\]
对递推变换，
\[
\mathrm{Val}_{\mathrm{pr}}(\pi_{k+2})
=F_{k+3}
=F_{k+2}+F_{k+1}
=\mathrm{Val}_{\mathrm{pr}}(\pi_{k+1}\pi_k).
\]
因此每一步局部替换都保持 \(\mathrm{Val}_{\mathrm{pr}}\) 不变，有限复合后仍保持不变。证毕。
\leanverified{base\_move\_preserves\_valPr}
\leanverified{fib\_move\_preserves\_valPr}
\leanverified{paper\_prime\_register\_local\_moves\_preserve\_valuation}
\end{proof}

\begin{theorem}[寄存器轨道与值纤维的一致性]\label{thm:prime-register-orbit-fiber-coincidence}
\leanverified{paper\_prime\_register\_orbit\_fiber\_coincidence}
对任意寄存器状态 \(r\in\mathcal R_{\Pi}\)，其 \(\mathcal G_{\Pi}\)-轨道恰等于其 Fibonacci 估值纤维：
\[
\mathcal G_{\Pi}\cdot r
=
\{s\in\mathcal R_{\Pi}:\ \mathrm{Val}_{\mathrm{pr}}(s)=\mathrm{Val}_{\mathrm{pr}}(r)\}.
\]
因此，两个寄存器状态可由局部 Fibonacci 变换互相到达，当且仅当它们具有相同的对象层值。
\end{theorem}

\begin{proof}
由命题 \ref{prop:prime-register-local-moves-preserve-valuation}，\(\mathcal G_{\Pi}\)-轨道必包含于值纤维。

反过来，若 \(\mathrm{Val}_{\mathrm{pr}}(r)=\mathrm{Val}_{\mathrm{pr}}(s)\)，记
\[
a=\Gamma^{-1}(r),
\qquad
b=\Gamma^{-1}(s).
\]
则 \(\mathrm{Val}(a)=\mathrm{Val}(b)\)。
由定理 \ref{thm:monoid-quotient-is-N} 与定理 \ref{thm:zeckendorf-transversal}，配置 \(a,b\) 属于同一 Fibonacci 同余类，并可由一系列局部可逆变换互达。
经由 \(\Gamma\) 转运，这些变换正对应 \(\mathcal R_{\Pi}\) 上的生成元，故 \(r,s\) 落在同一 \(\mathcal G_{\Pi}\)-轨道。证毕。
\end{proof}

\begin{definition}[折叠正规形寄存器与内部加法]\label{def:prime-register-normal-form}
对任意 \(r\in\mathcal R_{\Pi}\)，定义其折叠正规形为
\[
\mathrm{NF}_{\mathrm{pr}}(r):=R_F\bigl(\mathrm{Val}_{\mathrm{pr}}(r)\bigr)=\mathsf g\bigl(\mathrm{Val}_{\mathrm{pr}}(r)\bigr)\in\mathrm{ZG}.
\]
在 \(\mathrm{ZG}\) 上定义二元运算
\[
a\boxplus_{\mathrm{pr}} b:=\mathrm{NF}_{\mathrm{pr}}(ab)
\qquad(a,b\in\mathrm{ZG}).
\]
\end{definition}

\begin{proposition}[寄存器正规形的存在唯一性]\label{prop:prime-register-normal-form-uniqueness}
\leanverified{paper\_prime\_register\_normal\_form\_uniqueness}
对任意 \(r\in\mathcal R_{\Pi}\)，存在唯一折叠正规形寄存器 \(\mathrm{NF}_{\mathrm{pr}}(r)\in \mathrm{ZG}\)，并且 \(r\) 可经有限次局部值保持变换化到该状态。
\end{proposition}

\begin{proof}
经 \(\Gamma^{-1}\) 将 \(r\) 送回尺度计数配置空间，可得某个 \(a\in\mathcal D\)。
由命题 \ref{prop:directed-rewrite-termination}，\(a\) 经有限次值保持重写必终止于唯一 Zeckendorf 正规形 \(Z(\mathrm{Val}(a))\)。
再经 \(\Gamma\) 送回寄存器空间，即得到唯一寄存器正规形
\[
R_F(\mathrm{Val}(a))=R_F(\mathrm{Val}_{\mathrm{pr}}(r))=\mathrm{NF}_{\mathrm{pr}}(r).
\]
证毕。
\end{proof}

\begin{proposition}[寄存器替换重写系的终止性、合流性与正规形判定]\label{prop:prime-register-rewrite-confluent-normal-form}
在 \(\mathcal R_{\Pi}\) 上考虑如下局部值保持重写：
\begin{enumerate}
  \item 若 \(\pi_k\pi_{k+1}\mid r\)，则
  \[
  r\ \longrightarrow\ r\cdot \frac{\pi_{k+2}}{\pi_k\pi_{k+1}}
  \qquad(k\ge1);
  \]
  \item 若 \(\pi_1^2\mid r\)，则
  \[
  r\ \longrightarrow\ r\cdot \frac{\pi_2}{\pi_1^2};
  \]
  \item 若 \(\pi_2^2\mid r\)，则
  \[
  r\ \longrightarrow\ r\cdot \frac{\pi_1\pi_3}{\pi_2^2};
  \]
  \item 若 \(k\ge3\) 且 \(\pi_k^2\mid r\)，则
  \[
  r\ \longrightarrow\ r\cdot \frac{\pi_{k-2}\pi_{k+1}}{\pi_k^2}.
  \]
\end{enumerate}
则：
\begin{enumerate}
  \item 每一步重写都严格保持 \(\mathrm{Val}_{\mathrm{pr}}\) 不变；
  \item 该重写系在 \(\mathcal R_{\Pi}\) 上强终止，并具有唯一正规形；
  \item 对任意 \(r\in\mathcal R_{\Pi}\)，其唯一正规形恰为 \(\mathrm{NF}_{\mathrm{pr}}(r)\in\mathrm{ZG}\)。
\end{enumerate}
\end{proposition}

\begin{proof}
每条重写在指数向量层都对应 \S\ref{subsec:directed-rewrite} 中的局部规则：相邻合并、底位去重、第二位去重与一般位去重。
因此第 (1) 条由相应的 Fibonacci 权重恒等式
\[
F_{k+1}+F_{k+2}=F_{k+3},
\qquad
2F_2=F_3,
\qquad
2F_3=F_2+F_4,
\qquad
2F_{k+1}=F_{k-1}+F_{k+2}
\]
直接得到。

对第 (2)--(3) 条，记 \(\nu(r):=\Gamma^{-1}(r)\in\mathcal D\)。
上述整数重写在指数向量层恰对应命题 \ref{prop:directed-rewrite-termination} 的定向化归一化重写。
因此该重写系强终止，且终止形唯一，并等于
\[
Z(\mathrm{Val}(\nu(r)))=Z(\mathrm{Val}_{\mathrm{pr}}(r)).
\]
再经 \(\Gamma\) 送回寄存器空间，即得到唯一正规形
\[
\Gamma\bigl(Z(\mathrm{Val}_{\mathrm{pr}}(r))\bigr)
=R_F(\mathrm{Val}_{\mathrm{pr}}(r))
=\mathrm{NF}_{\mathrm{pr}}(r).
\]
证毕。
\end{proof}
\leanverified{paper\_prime\_register\_rewrite\_confluent\_normal\_form}

\begin{theorem}[乘法归一化实现可加结构的严格同构]\label{thm:prime-register-multiplicative-normalization-additive-iso}
\((\mathrm{ZG},\boxplus_{\mathrm{pr}})\) 为以 \(1\) 为单位元的交换幺半群，并且映射
\[
\mathsf g:\ (\NN,+)\longrightarrow (\mathrm{ZG},\boxplus_{\mathrm{pr}})
\]
是交换幺半群同构。
等价地，对任意 \(m,n\in\NN\)，有
\[
\mathrm{Val}_{\mathrm{pr}}(R_F(m)R_F(n))=m+n,
\qquad
\mathrm{NF}_{\mathrm{pr}}\bigl(R_F(m)R_F(n)\bigr)=R_F(m+n).
\]
\end{theorem}

\begin{proof}
首先，由定义 \ref{def:prime-register-normal-form} 与命题 \ref{prop:prime-register-internal-product-adds-values}，
\[
\mathrm{NF}_{\mathrm{pr}}\bigl(R_F(m)R_F(n)\bigr)
=R_F\!\Bigl(\mathrm{Val}_{\mathrm{pr}}(R_F(m)R_F(n))\Bigr)
=R_F(m+n).
\]
因此
\[
R_F(m)\boxplus_{\mathrm{pr}}R_F(n)=R_F(m+n),
\]
从而 \(\mathsf g=R_F\) 为幺半群同态。

其单射性由 Zeckendorf 规范形唯一性推出：若 \(R_F(m)=R_F(n)\)，则相应数字向量相同，故 \(m=n\)。
满射性由 \(\mathrm{ZG}=R_F(\NN)\) 的定义成立，故 \(\mathsf g\) 为交换幺半群同构。证毕。
\end{proof}
\leanverified{paper_prime_register_multiplicative_normalization_additive_iso}

\begin{remark}[加法的内部化与乘法的缺席]\label{rem:prime-register-addition-internalized}
定理 \ref{thm:prime-register-multiplicative-normalization-additive-iso} 表明，Fibonacci 寄存器已经把值的加法完全内部化：寄存器相乘，再经正规化，即对应对象层值相加。
然而对象层值的乘法并未在同一内部机制中出现。
若要实现乘法，必须把值先送回普通整数，再把其素因子指数外置到账本中。
这正是下文外置账本层的必要性。
\end{remark}

\begin{definition}[普通素数账本]\label{def:ordinary-prime-ledger}
记 \(\PP\) 为普通算术意义下的素数集合。
对每个正整数 \(n\ge1\)，定义其素数账本为
\[
\mathcal L(n)
:=
\sum_{p\in\PP}v_p(n)\,\mathbf e_p
\in
\bigoplus_{p\in\PP}\NN \mathbf e_p,
\]
其中 \(v_p(n)\) 为 \(n\) 的 \(p\)-进指数。

这里必须区分两套索引：
\(\Pi=(\pi_k)\) 记录的是 Fibonacci 尺度配置，而 \(\PP\) 记录的是对象层值本身的普通素因子分解。
前者适配加法归一化，后者适配乘法线性化。
\end{definition}

\begin{proposition}[普通素数账本线性化值的乘法]\label{prop:ordinary-prime-ledger-linearizes-multiplication}
对任意正整数 \(m,n\ge1\)，有
\[
\mathcal L(mn)=\mathcal L(m)+\mathcal L(n),
\qquad
\mathcal L(1)=0.
\]
并且 \(\mathcal L\) 为单射。
\leanverified{paper\_prime\_ledger\_linearizes\_multiplication}
\leanverified{paper\_ordinary\_prime\_ledger\_linearizes\_multiplication}
\end{proposition}

\begin{proof}
由唯一素因子分解，
\[
m=\prod_{p\in\PP}p^{v_p(m)},
\qquad
n=\prod_{p\in\PP}p^{v_p(n)},
\]
故
\[
mn=\prod_{p\in\PP}p^{v_p(m)+v_p(n)}.
\]
于是对每个 \(p\in\PP\) 有 \(v_p(mn)=v_p(m)+v_p(n)\)，从而
\[
\mathcal L(mn)
=\sum_{p\in\PP}(v_p(m)+v_p(n))\mathbf e_p
=\mathcal L(m)+\mathcal L(n).
\]
\(\mathcal L(1)=0\) 显然。
单射性同样来自唯一素因子分解：若 \(\mathcal L(m)=\mathcal L(n)\)，则一切 \(p\)-进指数相同，从而 \(m=n\)。证毕。
\end{proof}

\begin{theorem}[丢失信息的自由阿贝尔账本：值核与残差分数群]\label{thm:prime-register-residual-ledger-group}
令 \(A:=\bigoplus_{k\ge1}\ZZ e_k\) 为有限支撑整数序列的自由阿贝尔群，并定义群同态
\[
\varphi:A\to\ZZ,
\qquad
\varphi(e_k)=F_{k+1}.
\]
记 \(K:=\ker\varphi\)。
则：
\begin{enumerate}
  \item 记
  \[
  b_0:=e_2-2e_1,
  \qquad
  b_k:=e_{k+2}-e_{k+1}-e_k\quad(k\ge1).
  \]
  则 \(\{b_0,b_1,b_2,\dots\}\) 构成 \(K\) 的一组 \(\ZZ\)-基；特别地，\(K\) 为可数秩自由阿贝尔群。

  \item 定义素数分数群
  \[
  H:=
  \left\{
  \prod_{k\ge1}\pi_k^{\delta_k}
  \;:\;
  \delta=(\delta_k)_{k\ge1}\in K
  \right\}\subset\QQ_{>0}.
  \]
  则 \(H\) 亦为可数秩自由阿贝尔群，并由
  \[
  h_0:=\frac{\pi_2}{\pi_1^2},
  \qquad
  h_k:=\frac{\pi_{k+2}}{\pi_{k+1}\pi_k}\quad(k\ge1)
  \]
  自由生成。

  \item 对任意 \(r\in\mathcal R_{\Pi}\)，定义其残差分数
  \[
  R(r):=\frac{r}{\mathrm{NF}_{\mathrm{pr}}(r)}\in\QQ_{>0}.
  \]
  则 \(R(r)\in H\)，并且映射
  \[
  \Theta:\mathcal R_{\Pi}\to \mathrm{ZG}\times H,
  \qquad
  \Theta(r):=\bigl(\mathrm{NF}_{\mathrm{pr}}(r),R(r)\bigr)
  \]
  为单射。
  此外，\(\Gamma^{-1}(r)-Z(\mathrm{Val}_{\mathrm{pr}}(r))\in K\) 在基 \(\{b_i\}_{i\ge0}\) 上的坐标唯一存在，从而给出一份与归一化路径无关的 carry 账本。
\end{enumerate}
\leanverified{paper\_prime\_register\_residual\_ledger\_group}
\end{theorem}

\begin{proof}
\textbf{(1)} 由
\[
\varphi(b_0)=F_3-2F_2=0,
\qquad
\varphi(b_k)=F_{k+3}-F_{k+2}-F_{k+1}=0\quad(k\ge1),
\]
可知 \(b_i\in K\)。
对任意 \(N\ge2\)，令 \(A_N:=\bigoplus_{k=1}^{N}\ZZ e_k\)，\(K_N:=K\cap A_N\)。
则 \(\mathrm{rank}(K_N)=N-1\)。
集合 \(\{b_0,\dots,b_{N-2}\}\subset K_N\) 含有 \(N-1\) 个元素。
若 \(\sum_{i=0}^{N-2}c_i b_i=0\)，则从 \(e_N\) 的系数开始向下回代可得 \(c_{N-2}=\cdots=c_0=0\)，故它们线性无关，从而构成 \(K_N\) 的一组基。
由 \(K=\bigcup_{N}K_N\) 即得 \(\{b_i\}_{i\ge0}\) 为 \(K\) 的 \(\ZZ\)-基。

\textbf{(2)} 由算术基本定理，映射
\[
(\delta_k)_{k\ge1}\longmapsto \prod_{k\ge1}\pi_k^{\delta_k}
\]
给出 \(A\) 到由 \(\Pi\) 生成的分数群的群同构。
因此 \(H\simeq K\) 为可数秩自由阿贝尔群。
并且
\[
\prod_{k\ge1}\pi_k^{b_0}=\frac{\pi_2}{\pi_1^2}=h_0,
\qquad
\prod_{k\ge1}\pi_k^{b_k}=\frac{\pi_{k+2}}{\pi_{k+1}\pi_k}=h_k,
\]
故 \(\{h_i\}_{i\ge0}\) 为 \(H\) 的自由生成元。

\textbf{(3)} 设 \(r=\prod_{k\ge1}\pi_k^{a_k}\)，并记
\[
\mathrm{NF}_{\mathrm{pr}}(r)=\prod_{k\ge1}\pi_k^{z_k(\mathrm{Val}_{\mathrm{pr}}(r))}.
\]
于是
\[
R(r)=\prod_{k\ge1}\pi_k^{\delta_k(r)},
\qquad
\delta_k(r):=a_k-z_k(\mathrm{Val}_{\mathrm{pr}}(r)).
\]
由命题 \ref{prop:prime-register-rewrite-confluent-normal-form}，
\[
\mathrm{Val}_{\mathrm{pr}}(\mathrm{NF}_{\mathrm{pr}}(r))=\mathrm{Val}_{\mathrm{pr}}(r),
\]
故
\[
\sum_{k\ge1}\delta_k(r)F_{k+1}
=\mathrm{Val}_{\mathrm{pr}}(r)-\mathrm{Val}_{\mathrm{pr}}(\mathrm{NF}_{\mathrm{pr}}(r))
=0.
\]
因此 \(\delta(r):=(\delta_k(r))\in K\)，从而 \(R(r)\in H\)。

单射性由指数向量的逐坐标恢复给出：\(\Theta(r)=(a,h)\) 等价于
\[
\Gamma^{-1}(r)=\Gamma^{-1}(a)+\delta(h),
\]
其中 \(\delta(h)\in K\subset A\) 为 \(h\) 的唯一指数向量。
右端唯一确定 \(\Gamma^{-1}(r)\)，继而唯一确定 \(r\)。
最后，\(\Gamma^{-1}(r)-Z(\mathrm{Val}_{\mathrm{pr}}(r))\in K\) 在基 \(\{b_i\}\) 上坐标唯一存在，且仅由 \(r\) 与其正规形决定，故与任何具体重写路径无关。证毕。
\end{proof}

\begin{definition}[Fibonacci 值的外置账本]\label{def:prime-register-external-ledger}
对任意满足 \(\mathrm{Val}_{\mathrm{pr}}(r)\ge1\) 的寄存器状态 \(r\in\mathcal R_{\Pi}\)，定义其外置账本为
\[
\Lambda_F(r):=\mathcal L(\mathrm{Val}_{\mathrm{pr}}(r)).
\]
也就是说，先从 Fibonacci 寄存器中恢复对象层值，再把该值的乘法结构送入普通素数账本。
\end{definition}

\begin{proposition}[外置账本在值保持轨道上不变]\label{prop:prime-register-external-ledger-orbit-invariance}
若 \(r,s\in\mathcal R_{\Pi}\) 满足 \(s\in\mathcal G_{\Pi}\cdot r\)，且 \(\mathrm{Val}_{\mathrm{pr}}(r)\ge1\)，则
\[
\Lambda_F(r)=\Lambda_F(s).
\]
特别地，
\[
\Lambda_F(r)=\Lambda_F(\mathrm{NF}_{\mathrm{pr}}(r)).
\]
\leanverified{paper\_prime\_register\_external\_ledger\_orbit\_invariance}
\end{proposition}

\begin{proof}
由定理 \ref{thm:prime-register-orbit-fiber-coincidence}，\(s\in\mathcal G_{\Pi}\cdot r\) 等价于 \(\mathrm{Val}_{\mathrm{pr}}(s)=\mathrm{Val}_{\mathrm{pr}}(r)\)。
因此
\[
\Lambda_F(s)=\mathcal L(\mathrm{Val}_{\mathrm{pr}}(s))
=\mathcal L(\mathrm{Val}_{\mathrm{pr}}(r))
=\Lambda_F(r).
\]
取 \(s=\mathrm{NF}_{\mathrm{pr}}(r)\) 即得第二式。证毕。
\end{proof}

\begin{definition}[拉回的值乘]\label{def:prime-register-pulled-back-value-multiplication}
在规范寄存器集合
\[
\mathcal N_F:=\{R_F(n):n\in\NN\}=\mathrm{ZG}
\]
上定义新运算 \(\odot\)：
\[
R_F(m)\odot R_F(n):=R_F(mn).
\]
这不是寄存器空间原有的整数乘法，而是把对象层整数乘法经由 \(R_F\) 拉回到规范寄存器集合上的结果。
\end{definition}

\begin{proposition}[内部乘法与拉回值乘的严格区分]\label{prop:prime-register-internal-vs-pulled-back-product}
\leanverified{paper\_prime\_register\_internal\_vs\_pulled\_back\_product}
对任意 \(m,n\in\NN_{>0}\)，有
\[
R_F(m)R_F(n)\xrightarrow{\ \mathrm{NF}_{\mathrm{pr}}\ }R_F(m+n),
\]
而
\[
R_F(m)\odot R_F(n)=R_F(mn).
\]
此外，外置账本对 \(\odot\) 线性化：
\[
\Lambda_F\bigl(R_F(m)\odot R_F(n)\bigr)
=\Lambda_F(R_F(m))+\Lambda_F(R_F(n)).
\]
\end{proposition}

\begin{proof}
第一式即定理 \ref{thm:prime-register-multiplicative-normalization-additive-iso}。
第二式由定义 \ref{def:prime-register-pulled-back-value-multiplication}。
第三式由定义 \ref{def:prime-register-external-ledger} 与命题 \ref{prop:ordinary-prime-ledger-linearizes-multiplication}：
\[
\Lambda_F(R_F(m)\odot R_F(n))
=\mathcal L(mn)
=\mathcal L(m)+\mathcal L(n)
=\Lambda_F(R_F(m))+\Lambda_F(R_F(n)).
\]
证毕。
\end{proof}

\begin{definition}[有限可见账本与余额余因子]\label{def:prime-register-finite-ledger}
对有限素数集 \(S\subset\PP\)，定义截断账本
\[
\mathcal L_S(n):=\sum_{p\in S}v_p(n)\,\mathbf e_p,
\]
并定义余额余因子
\[
\rho_S(n):=
n\cdot\prod_{p\in S}p^{-v_p(n)}.
\]
于是 \(n\) 可唯一写成
\[
n=\left(\prod_{p\in S}p^{v_p(n)}\right)\rho_S(n),
\]
其中前半部分对应账本窗口 \(S\) 可见的素数分量，后半部分则是窗口之外的余额。
\end{definition}

\begin{proposition}[有限账本的可恢复性]\label{prop:prime-register-finite-ledger-recoverability}
对任意有限素数集 \(S\subset\PP\)，映射
\[
n\longmapsto \bigl(\mathcal L_S(n),\rho_S(n)\bigr)
\]
是单射。
因此，有限截断账本本身不足以恢复对象层值，但“截断账本 + 余额余因子”能够完全恢复原值。
\end{proposition}

\begin{proof}
由定义 \ref{def:prime-register-finite-ledger}，
\[
n=
\left(\prod_{p\in S}p^{(\mathcal L_S(n))_p}\right)\rho_S(n).
\]
右端由二元组 \((\mathcal L_S(n),\rho_S(n))\) 唯一确定，故该映射为单射。证毕。
\end{proof}
\leanverified{paper\_prime\_register\_finite\_ledger\_recoverability}

\begin{remark}[与对象层 forcing 语义的衔接]\label{rem:prime-register-forcing-interface}
本小节的一切构造都以“值已经进入对象层”为前提。
只有当某个地址--值对在第 \ref{sec:recursive-addressing} 节的意义下已经完成对象化、通过协议允许性与非 \(\mathsf{NULL}\) 判据时，才谈得上为其建立 \(R_F(n)\)、\(\mathrm{NF}_{\mathrm{pr}}\)、\(\Lambda_F\) 或 \(\mathcal L_S\) 等值层结构。
因此，本小节并不替代对象层准入，而是在其基础上继续做值层实现：第 \ref{sec:recursive-addressing} 节回答何时可谈，第 \ref{subsec:emergent-arithmetic-prime-register-fib-valuation} 节回答既然已可谈，如何把该值实现为 Fibonacci 寄存器正规形，并把其乘法结构外置到账本。
\end{remark}

\begin{remark}[本小节的两层结构]\label{rem:prime-register-two-layer-summary}
至此得到两层互补而不可混同的结构。
第一层是 Fibonacci 寄存器层：它把尺度计数配置实现为素数指数向量，把值映射实现为 Fibonacci 加权估值，并把规范化实现为局部值保持重写通向唯一折叠正规形；因此，这一层天然适配加法与归一化。

第二层是普通素数账本层：它不再记录尺度配置，而记录对象层值本身的素因子指数；因而它把值的乘法线性化为加法，但其位置必须是外置的。
有限账本只给出局部乘法可见性，而完整恢复还需余额余因子。
\end{remark}

\subsubsection{动态 prime-register：Gödel 外置动力学}\label{subsubsec:emergent-arithmetic-dynamic-prime-register-godel-dynamics}

前两小节已经把 prime-register 的静态骨架固定下来。
一方面，推论 \ref{cor:killo-no-finite-additive-register-linearization} 与定理 \ref{thm:killo-godel-compression-not-finite-rank-homologizable} 表明：若坚持“乘法 \(\to\) 寄存器加法”的严格同态语义，则有限个加法寄存器不足以忠实承载 \((\NN_{\ge 1},\times)\)，而可数无穷素数轴
\[
n\longmapsto (v_p(n))_{p\in\PP}
\]
给出最小可行的严格同态骨架。
另一方面，定理 \ref{thm:killo-localization-circle-dimension-prime-ledger-splitting} 又把局部化数系的对偶分解写成
\[
0\to \prod_{p\in S}\ZZ_p\to \widehat{G_S}\to \TT\to 0,
\]
从而把连通的相位圆维与全不连通的素数账本核严格区分开，并证明有限支持集合 \(S\) 可由该核完全恢复。
这说明 prime-register 在本文中并不是附加编码技巧，而是乘法保持、可逆外置与历史可审计同时成立时被结构强迫出来的最小账本。

但到目前为止，这一账本主要以静态外置出现：它告诉我们对象层值如何被编码，却尚未单独对象化“编码如何随时间推进”。
本小节的任务就是补上这一步。
其核心思想是：Gödel 编码不应只被看成单个对象的标签，而应被提升为事件词、前史地址与时间顺序的动力学账本。
一旦这样做，prime-register 就不再只是乘法层坐标系，而会成为一条真正的时间外置轴；地址精度提升、前缀延长、事件流拼接与几何实现，都可在同一半直积框架中统一书写。

\begin{definition}[动态 prime-register 过程]\label{def:emergent-arithmetic-dynamic-prime-register-process}
设 \(\mathcal E\) 为原子事件集合，\(\mathrm{code}:\mathcal E\hookrightarrow \NN_{\ge 1}\) 为单射编码。
记素数序列为
\[
p_1<p_2<\cdots,
\qquad
\mathcal P:=\NN^{(\PP)}.
\]
对任意事件流 \((e_t)_{t\ge 1}\in \mathcal E^{\NN}\) 及其长度为 \(T\) 的前缀，定义动态 prime-register
\[
r_T(p_t):=
\begin{cases}
\mathrm{code}(e_t),& 1\le t\le T,\\
0,& t>T,
\end{cases}
\]
以及对应的 Gödel 前缀整数
\[
N_T:=N(r_T)=\prod_{t=1}^{T}p_t^{\,\mathrm{code}(e_t)}.
\]
若事件流来自某个输入 \(\omega\) 的归一化、回放或自然扩张策略，则称 \((r_T(\omega),N_T(\omega))\) 为该输入的动态 Gödel 账本前缀。
\end{definition}

这个定义把静态 prime-register 外置提升为真正的时间对象：\(T\) 不再只是外部时钟，而是账本前缀长度；\(r_T\) 不再只是指数向量，而是带有时间寻址的前史地址；\(N_T\) 则是该前史地址在交换乘法层上的整数影子。

\begin{theorem}[动态账本的拼接律与时间寻址]\label{thm:emergent-arithmetic-dynamic-prime-register-concatenation}
\leanverified{paper\_emergent\_arithmetic\_dynamic\_prime\_register\_concatenation}
设 \(\Sigma^k:\mathcal P\to \mathcal P\) 为移位端同态
\[
(\Sigma^k r)(p_t):=
\begin{cases}
0,& t\le k,\\
r(p_{t-k}),& t>k,
\end{cases}
\]
并定义半直积幺半群
\[
\mathcal H:=\mathcal P\rtimes_{\Sigma}\NN,
\qquad
(r,k)\cdot(s,\ell):=(r+\Sigma^k s,\ k+\ell).
\]
对任意有限事件词 \(w=e_1\cdots e_T\in\mathcal E^\ast\)，记
\[
r_w(p_t):=
\begin{cases}
\mathrm{code}(e_t),& 1\le t\le T,\\
0,& t>T,
\end{cases}
\qquad
h(w):=(r_w,T)\in \mathcal H.
\]
则对任意 \(u,v\in\mathcal E^\ast\)，有
\[
h(uv)=h(u)\cdot h(v),
\]
亦即
\[
r_{uv}=r_u+\Sigma^{|u|}r_v,
\qquad
|uv|=|u|+|v|.
\]
等价地，若记 \(N_w:=N(r_w)\)，则存在由 \(p_t\mapsto p_{t+k}\) 诱导的整数端同态 \(\mathrm{Sh}_k\) 使
\[
N_{uv}=N_u\cdot \mathrm{Sh}_{|u|}(N_v).
\]
\end{theorem}

\begin{proof}
这是定理 \ref{thm:conclusion-godel-semidirect-law} 的直接转写。
该定理已经证明 \(h:\mathcal E^\ast\to \mathcal H\) 为幺半群嵌入，并满足
\[
h(uv)=h(u)\cdot h(v).
\]
把半直积乘法
\[
(r,k)\cdot(s,\ell)=(r+\Sigma^k s,\ k+\ell)
\]
展开，即得
\[
r_{uv}=r_u+\Sigma^{|u|}r_v,
\qquad
|uv|=|u|+|v|.
\]
再由同一结果中的整数化形式
\[
N( r+s)=N(r)N(s),
\qquad
N(\Sigma^k r)=\mathrm{Sh}_k(N(r))
\]
便得
\[
N_{uv}=N_u\cdot \mathrm{Sh}_{|u|}(N_v).
\]
证毕。
\end{proof}

\begin{remark}[顺序不是附加注释，而是 prime-slice 位移]\label{rem:emergent-arithmetic-dynamic-prime-register-time-addressing}
定理 \ref{thm:emergent-arithmetic-dynamic-prime-register-concatenation} 的含义在于：动态 Gödel 账本并不是把事件词简单压缩成一个大整数，而是把“连接 = 位移后相加”的时间拼接律直接内建到 prime-slices 中。
因此，顺序信息不是靠外部时间标签补上，而是由 \(r_u+\Sigma^{|u|}r_v\) 这一位移规则内生保留。
这也解释了为何普通单整数乘法同态不足以承载非交换顺序：一旦去掉时间位移，剩下的就只是一条交换乘法通道。
\end{remark}

\begin{proposition}[有限窗口账本的可行域]\label{prop:emergent-arithmetic-dynamic-prime-register-bounded-feasibility}
\leanverified{paper\_emergent\_arithmetic\_dynamic\_prime\_register\_bounded\_feasibility}
设 \(f:\Omega\twoheadrightarrow X\) 为有限集合间满射，并记
\[
D(f):=\max_{x\in X}|f^{-1}(x)|.
\]
对 \(k,E\in\NN\)，令
\[
\mathcal P_{k,E}
:=
\left\{
r\in\mathcal P:\ r(p_i)\in\{0,\dots,E\}\ (1\le i\le k),\ r(p)=0\ (p>p_k)
\right\}.
\]
则存在单射
\[
\iota:\Omega\to X\times\mathcal P_{k,E},
\qquad
\operatorname{pr}_X\circ\iota=f,
\]
当且仅当
\[
(E+1)^k\ge D(f).
\]
特别地，对 \(f=\Fold_m\) 而言，若
\[
\iota_m:\Omega_m\to X_m\times\mathcal P_{k,E},
\qquad
\operatorname{pr}_X\circ\iota_m=\Fold_m,
\]
为单射，则必有
\[
(E+1)^k\ge D_m,
\qquad
D_m:=\max_{x\in X_m}|\Fold_m^{-1}(x)|.
\]
\end{proposition}

\begin{proof}
前半部分正是定理 \ref{thm:conclusion-bounded-prime-register-feasibility} 的内容；后半部分则是将该定理应用于 \(f=\Fold_m\) 的直接结论。
进一步，对 \(\Fold_m\) 的最大纤维尺度，定理 \ref{thm:conclusion-fold-prime-axis-exponent-lower-bound} 已给出
\[
D_{2j}=F_{j+2},
\qquad
D_{2j+1}=2F_{j+1},
\]
从而
\[
\log D_m=\frac{m}{2}\log\varphi+O(1).
\]
这说明有限窗口下的忠实外置虽可由有限 prime-slices 实现，但轴数与指数上界之间必须满足严格的乘法容量约束。证毕。
\end{proof}

\begin{remark}[全深度前史强迫无限 prime-support]\label{rem:emergent-arithmetic-dynamic-prime-register-infinite-support}
有限窗口的可行性并不意味着全深度前史也能被固定有限账本吸收。
相反，推论 \ref{cor:conclusion-faithful-time-addressed-godel-needs-infinite-prime-support} 已表明：任何 shift-compatible 且 faithful 的时间寻址 Gödel 编码都必然使用无限素数支撑。
若进一步要求与自然扩张的真实分支历史兼容，则推论 \ref{cor:xi-time-part9zq-natural-extension-infinite-primeslice-support} 给出更强的结论：只要每一深度都存在非平凡分支，则任一与有限截断相容的 injective prime-slice 账本在无限深度上都必须具有无限活跃片支撑。
于是，定理 \ref{thm:xi-time-part9zq-no-fixed-finite-support-solenoid-host} 进一步排除了把全部无穷分支历史压入某个固定有限支持局部化系统 \(G_S\) 或其 Pontryagin 对偶 \(\widehat{G_S}\) 的可能性。
因此，动态 prime-register 必须严格区分两层问题：有限窗口的可审计外置与无穷前史的忠实宿主，一般并不是同一尺度上的对象。
\end{remark}

\begin{proposition}[动态 Gödel 账本的位长代价]\label{prop:emergent-arithmetic-dynamic-prime-register-bitlength}
沿用定义 \ref{def:pom-fold-prime-lift} 中的策略 \(\mathsf S\) 与 prime-register 回放记号。
对输入 \(\omega\in\Omega_{m+1}\)，记
\[
G_{\mathsf S}(\omega):=N\!\bigl(r_{\mathsf S}(\omega)\bigr)
=\prod_{t=1}^{T_{m+1}^{\mathsf S}(\omega)}
p_t^{\,\mathrm{code}(e_t(\omega))}.
\]
则对任意 \(\omega\) 都有点态下界
\[
\log_2 G_{\mathsf S}(\omega)
\ge
\log_2\!\bigl((T_{m+1}^{\mathsf S}(\omega)+1)!\bigr),
\]
从而
\[
\log_2 G_{\mathsf S}(\omega)
=
\Omega\!\bigl(T_{m+1}^{\mathsf S}(\omega)\log T_{m+1}^{\mathsf S}(\omega)\bigr).
\]
并且在均匀基线 \(\omega\sim\mathrm{Unif}(\Omega_{m+1})\) 下，存在常数 \(C>0\) 使
\[
\EE[\log_2 G_{\mathsf S}(\omega)]
\ge
\frac{4}{27\ln 2}\,m\log m-Cm.
\]
\leanverified{paper\_emergent\_arithmetic\_dynamic\_prime\_register\_bitlength}
\end{proposition}

\begin{proof}
这是定理 \ref{thm:conclusion-prime-integerization-superlinear-bitlength} 的直接转写；同样的 \(m\log m\) 下界也由定理 \ref{thm:pom-fold-prime-godel-bitlength-superlinear} 给出。
其证明机制很简单：由 \(\mathrm{code}(e_t)\ge 1\) 可知
\[
G_{\mathsf S}(\omega)\ge \prod_{t=1}^{T_{m+1}^{\mathsf S}(\omega)}p_t\ge (T_{m+1}^{\mathsf S}(\omega)+1)!,
\]
再结合 Stirling 型下界与事件步数的线性期望下界，便得到点态 \(T\log T\) 量级以及均值 \(m\log m\) 量级。
证毕。
\end{proof}

\begin{remark}[位长增长就是时间外置成本]\label{rem:emergent-arithmetic-dynamic-prime-register-bitlength-cost}
命题 \ref{prop:emergent-arithmetic-dynamic-prime-register-bitlength} 的意义在于：动态账本一旦承担可回放前史，就不再是“免费可逆化”。
时间前缀越长，Gödel 整数位长就越昂贵；因此这里的时间不仅是步骤数，也是一种不可压缩的外置成本。
这正把 prime-register 账本与后文关于同步、资源与解析深度的讨论接通起来。
\end{remark}

\begin{theorem}[\texorpdfstring{$2^L$}{2\^L}-进仿射宿主中的动态账本]\label{thm:emergent-arithmetic-dynamic-prime-register-2adic-affine-host}
\leanverified{paper\_emergent\_arithmetic\_dynamic\_prime\_register\_2adic\_affine\_host}
设 \(\mathcal E\) 为有限原子事件集，并固定单射编码
\[
\mathrm{code}:\mathcal E\hookrightarrow \NN,
\qquad
M:=\max_{e\in\mathcal E}\mathrm{code}(e).
\]
取 \(B:=2^L>M\)，记
\[
T:=T_2(E_{\min})\cong \ZZ_2^2,
\]
并固定非零向量 \(v\in T\)。
定义
\[
\mathrm{ev}_B(r):=\sum_{t\ge 1}r(p_t)B^{t-1}v
\qquad (r\in\mathcal P),
\]
以及
\[
\Psi:\mathcal H\to \mathrm{Aff}(T),
\qquad
\Psi(r,k)(x):=B^k x+\mathrm{ev}_B(r).
\]
则：
\begin{enumerate}
  \item \(\Psi\) 是从半直积幺半群 \(\mathcal H=\mathcal P\rtimes_{\Sigma}\NN\) 到 \(\mathrm{Aff}(T)\) 的单子同态；
  \item 在规范编码子集上，\(\Psi\) 为忠实表示；
  \item 对任意无限事件流 \(\mathbf e=(e_1,e_2,\dots)\in\mathcal E^{\NN}\)，级数
  \[
  X(\mathbf e):=\sum_{t\ge 1}\mathrm{code}(e_t)B^{t-1}v
  \]
  在 \(2\)-进拓扑下收敛；其前缀和
  \[
  X_n(\mathbf e):=\sum_{t=1}^{n}\mathrm{code}(e_t)B^{t-1}v
  \]
  满足
  \[
  X_n(\mathbf e)\equiv X(\mathbf e)\pmod{B^nT},
  \]
  且映射 \(\mathbf e\mapsto X(\mathbf e)\) 为单射。
\end{enumerate}
\end{theorem}

\begin{proof}
第 (1)--(2) 条正是定理 \ref{thm:killo-godel-semidir-affine-2adic} 的内容；第 (3) 条则由推论 \ref{cor:killo-infinite-stream-2adic-holographic-point} 直接给出。
其中关键恒等式是
\[
\mathrm{ev}_B(\Sigma^k r)=B^k\mathrm{ev}_B(r),
\]
它把 prime-slice 位移精确转写为 \(2^L\)-进尺度放大。
因此
\[
\Psi(r,k)\circ \Psi(s,\ell)=\Psi(r+\Sigma^k s,\ k+\ell),
\]
而无限流的前缀和则在 \(B^nT\) 意义下逼近唯一的全息点 \(X(\mathbf e)\)。
证毕。
\end{proof}

\begin{remark}[时间推进可解释为地址精度提升]\label{rem:emergent-arithmetic-dynamic-prime-register-address-refinement}
定理 \ref{thm:emergent-arithmetic-dynamic-prime-register-2adic-affine-host} 给出了一种特别适合后文继续使用的解释：
在这个 \(2^L\)-进仿射宿主中，时间推进并不是“点在外部时间轴上运动”，而是同一个地址点被越来越精细地读取。
半直积中的 \(k\) 负责尺度提升，\(\mathrm{ev}_B(r)\) 负责平移注入，而 \(X_n(\mathbf e)\equiv X(\mathbf e)\pmod{B^nT}\) 则说明第 \(n\) 步前缀只是在同一全息地址上增加了一位可解析精度。
\end{remark}

\begin{remark}[与 prime-indexed 粗糙读出的接口]\label{rem:emergent-arithmetic-dynamic-prime-register-rough-readout-interface}
动态 prime-register 一旦被对象化，6.5.2 中的多尺度粗糙读出就获得了一种更贴合本文语义的频率来源：高频层不必再来自任意的 lacunary 频率族，而可以直接由 prime-slices 或 \(B\)-进地址层产生。
因此，最自然的下一步不是直接把 \(N_T\) 当作可见量，而是考虑由动态账本前缀调制相位的读出，例如
\[
W_\omega(t)
:=
\sum_{j\ge 1}
a_j
\cos\!\bigl(\lambda_j t+\phi_j(r_j(\omega))\bigr),
\]
其中 \(\lambda_j\) 可取 \(p_j\)、\(\log p_j\) 或 \(B^j\) 型尺度，而 \(\phi_j\) 则由第 \(j\) 个 prime-slice 上的账本前缀决定。
在这一接口下，6.5.2 的“确定性产生不可导可见量”、14.2.9 的“粗糙可见量在有限分辨率下表现为表观随机”，以及本节的“prime-register 是乘法与前史的最小动态账本”，便可被读成同一生成机制的不同投影面。
\end{remark}


在此基础上，外置尺度不再停留在离散指数账本层，而被进一步送入一个可解析观测的复几何门结构。
下一小节给出与该外置寄存器严格兼容的 Gödel--Joukowsky 椭圆化，并据此讨论解析不可见性、径向恢复与零点输运。

\subsubsection{Gödel--Joukowsky 椭圆化门：素数外置尺度的解析不可见性与零点输运}\label{subsec:prime-register-godel-joukowsky}

上一小节已经把 Fibonacci 寄存器的内部加法结构与对象层值的外置乘法账本严格分离。
以下固定定理~\ref{thm:prime-register-residual-ledger-group} 的第~(2)~部分所定义的残差分数群 \(H\)，并把这些外置尺度几何化为面积保持椭圆族。
这样，乘法侧的信息不再停留在离散指数账本层，而被送入一个可解析观测的复几何门结构；其解析可见性与不可见性，正是后续零点输运的出发点。

\begin{theorem}[值保持算法群的稠密椭圆化]\label{thm:prime-register-dense-ellipticization}
令 \(H\) 如定理~\ref{thm:prime-register-residual-ledger-group} 的第~(2)~部分。
考虑映射
\[
\Psi: H\to \mathrm{SL}(2,\RR),\qquad r\longmapsto
\begin{pmatrix}
r & 0\\[2pt]
0 & r^{-1}
\end{pmatrix}.
\]
并令 \(\mathcal E\) 为所有与坐标轴对齐且面积为 \(\pi\) 的椭圆集合，即
\[
\mathcal E=\{\Psi(r)(S^1): r>0\},
\qquad
S^1:=\{(x,y)\in\RR^2:\ x^2+y^2=1\}.
\]
则：
\begin{enumerate}
  \item \(\Psi(H)\) 在 \(\Psi(\RR_{>0})\) 中稠密；等价地，\(\{\Psi(r)(S^1):r\in H\}\) 在 \(\mathcal E\) 中稠密；
  \item 每个 \(r\in H\) 唯一决定椭圆 \(\Psi(r)(S^1)\)，且该椭圆的轴比为 \(r^2\)。
  因而 \(r\mapsto \Psi(r)(S^1)\) 给出一条把外置账本几何化为面积保持椭圆的可计算函子。
\end{enumerate}
\leanverified{paper\_prime\_register\_dense\_ellipticization}
\leanverified{diagAction}
\leanverified{diag\_maps\_circle\_to\_ellipse}
\leanverified{ellipse\_axes\_product}
\leanverified{axis\_ratio\_eq\_r\_sq}
\leanverified{r\_unique\_from\_sq}
\leanverified{ellipse\_area\_eq\_pi}
\end{theorem}

\begin{proof}
\textbf{(2)} 对角矩阵作用将单位圆送到半轴长度分别为 \(r\) 与 \(r^{-1}\) 的椭圆，面积为 \(\pi rr^{-1}=\pi\)，轴比为 \(r/r^{-1}=r^2\)，并由轴比唯一恢复 \(r>0\)。

\textbf{(1)} 只需证明 \(H\) 在 \(\RR_{>0}\) 中稠密。
取两元素
\[
a:=\frac{\pi_2}{\pi_1^2}\in H,
\qquad
b:=\frac{\pi_3}{\pi_1^3}\in H,
\]
其中 \(b\in H\) 对应核元 \((-3,0,1,0,\dots)\in K\)，因为 \(-3F_2+F_4=-3+3=0\)。
记 \(\alpha:=\log a\)、\(\beta:=\log b\)。
若 \(\alpha/\beta\in\QQ\)，则存在非零整数 \(m,n\) 使 \(m\alpha+n\beta=0\)，即 \(a^m b^n=1\)。
但
\[
a^m b^n
=\left(\frac{\pi_2}{\pi_1^2}\right)^m\left(\frac{\pi_3}{\pi_1^3}\right)^n
=\frac{\pi_2^m\pi_3^n}{\pi_1^{2m+3n}},
\]
由唯一素因子分解得 \(m=n=0\)，矛盾。
故 \(\alpha/\beta\notin\QQ\)。

于是加法子群
\[
L:=\{m\alpha+n\beta:\ m,n\in\ZZ\}\subset\RR
\]
为稠密子集。
指数化得到
\[
\exp(L)=\{a^m b^n:\ m,n\in\ZZ\}\subset H
\]
在 \(\RR_{>0}\) 中稠密，从而 \(H\) 稠密。
由 \(r\mapsto \Psi(r)\) 的连续性，\(\Psi(H)\) 在 \(\Psi(\RR_{>0})\) 中稠密；椭圆像的稠密性同理。证毕。
\end{proof}

本小节以下给出一条与外置素数指数口径直接兼容的复迹椭圆化：把乘法尺度 \(r\) 视为可审计外置量，并以 Joukowsky 型投影把单位圆 \(\TT\) 的相位信息输运到焦点固定的椭圆族 \(\{\mathcal E_r\}\subset\CC\) 上。
该椭圆族在解析观测下呈现刚性不可见性，而在径向非解析观测下保持可逆。

\begin{definition}[Gödel 尺度与 Joukowsky--Gödel 投影]\label{def:prime-register-godel-scale-joukowsky}
记 \(\ZZ^{(\PP)}\) 为仅有限支撑的素数指数向量群：
\[
\ZZ^{(\PP)}:=\Bigl\{a=(a_p)_{p\in\PP}:\ a_p\in\ZZ,\ a_p=0\text{ 除有限多素数外}\Bigr\}.
\]
定义 Gödel 尺度为
\begin{equation}\label{eq:prime-register-godel-scale}
\boxed{\;
\mathsf{G}(a):=\prod_{p\in\PP}p^{a_p}\in\QQ_{>0}\; }.
\end{equation}
并记 \(\NN^{(\PP)}:=\{a\in\ZZ^{(\PP)}:\ a_p\ge 0\ \forall p\}\) 为其非负子幺半群。

对任意 \(r>0\) 定义 Joukowsky--Gödel 投影
\begin{equation}\label{eq:prime-register-joukowsky-godel}
\boxed{\;
J_r:\ \TT\to\CC,\qquad
J_r(e^{\mathrm{i}\theta})
:=
r e^{\mathrm{i}\theta}+r^{-1}e^{-\mathrm{i}\theta}\; }.
\end{equation}
并记其像为
\[
\mathcal E_r:=J_r(\TT)\subset\CC.
\]
令 \(\mu_r\) 为 Haar 测度 \(\frac{d\theta}{2\pi}\) 在 \(J_r\) 下的推前概率测度。
\end{definition}

\begin{proposition}[Gödel--Joukowsky 椭圆化：素数寄存器决定椭圆]\label{prop:prime-register-godel-joukowsky-ellipse}
对任意 \(r>0\)，集合 \(\mathcal E_r\) 为平面椭圆（退化情形 \(r=1\) 为线段 \([-2,2]\)），并满足：
\begin{enumerate}
  \item \(\mathcal E_r=\mathcal E_{1/r}\)。在规范 \(r\ge1\) 下，椭圆 \(\mathcal E_r\) 唯一确定参数 \(r\)；
  \item \(\mathcal E_r\) 的焦点为 \(\pm2\)。在 \(r\ge1\) 下，其半长轴与半短轴分别为
  \[
  A(r)=r+r^{-1},
  \qquad
  B(r)=r-r^{-1};
  \]
  \item 在 \(r\ge1\) 的规范下，若定义
  \[
  \mathcal E_r\odot\mathcal E_s:=\mathcal E_{rs},
  \]
  则映射
  \[
  \Phi:\ \NN^{(\PP)}\to\{\mathcal E_r:\ r\ge1\},
  \qquad
  a\longmapsto \mathcal E_{\mathsf G(a)}
  \]
  为交换幺半群同态。
\end{enumerate}
\end{proposition}

\begin{proof}
由 \eqref{eq:prime-register-joukowsky-godel} 展开得
\[
J_r(e^{\mathrm{i}\theta})
=(r+r^{-1})\cos\theta+\mathrm{i}(r-r^{-1})\sin\theta,
\]
从而 \(\mathcal E_r\) 满足
\[
\left(\frac{\Re w}{r+r^{-1}}\right)^2+\left(\frac{\Im w}{|r-r^{-1}|}\right)^2=1.
\]
故其为椭圆；当 \(r=1\) 时退化为 \([-2,2]\)。
该方程仅依赖 \(\{r,r^{-1}\}\)，故 \(\mathcal E_r=\mathcal E_{1/r}\)。

当 \(r\ge1\) 时，
\[
A(r)^2-B(r)^2=(r+r^{-1})^2-(r-r^{-1})^2=4,
\]
故焦距为 \(4\)，焦点为 \(\pm2\)。
又有 \(A(r)+B(r)=2r\)，遂可由椭圆唯一恢复 \(r\)。

最后，由 \(\mathsf G(a+b)=\mathsf G(a)\mathsf G(b)\) 得
\[
\Phi(a+b)=\mathcal E_{\mathsf G(a+b)}
=\mathcal E_{\mathsf G(a)\mathsf G(b)}
=\Phi(a)\odot \Phi(b).
\]
证毕。
\end{proof}
\leanverified{paper\_prime\_register\_godel\_joukowsky\_ellipse}

\begin{theorem}[Cauchy 变换刚性：外置尺度对解析观测不可见]\label{thm:prime-register-joukowsky-cauchy-rigidity}
在定义 \ref{def:prime-register-godel-scale-joukowsky} 的记号下，对 \(z\in\CC\setminus\mathcal E_r\) 定义 Cauchy 变换
\[
G_r(z):=\int_{\mathcal E_r}\frac{1}{z-w}\,d\mu_r(w).
\]
取平方根分支使 \(\sqrt{z^2-4}\sim z\)（\(z\to\infty\)）。
则：
\begin{enumerate}
  \item 若 \(z\) 属于 \(\CC\setminus\mathcal E_r\) 的无界连通分支，则
  \[
  \boxed{\;G_r(z)=\frac{1}{\sqrt{z^2-4}}\;}
  \]
  因而在该解析观测下完全不依赖于 \(r\)；
  \item 若 \(r\ne1\) 且 \(z\) 位于 \(\mathcal E_r\) 围成的内部，则
  \[
  \boxed{\;G_r(z)=0.\;}
  \]
\end{enumerate}
\end{theorem}

\begin{proof}
由 \(J_{1/r}(e^{\mathrm{i}\theta})=J_r(e^{-\mathrm{i}\theta})\) 与 Haar 测度在 \(\theta\mapsto-\theta\) 下的不变性，可无损假设 \(r\ge1\)。

由推前测度定义与代换 \(\zeta=e^{\mathrm{i}\theta}\) 得
\[
G_r(z)=\frac{1}{2\pi\mathrm{i}}\oint_{|\zeta|=1}\frac{d\zeta}{-r\zeta^2+z\zeta-r^{-1}}.
\]
令
\[
D(\zeta):=-r\zeta^2+z\zeta-r^{-1},
\qquad
\zeta_{\pm}:=\frac{z\pm\sqrt{z^2-4}}{2r},
\]
则 \(D(\zeta)=-r(\zeta-\zeta_+)(\zeta-\zeta_-)\)，并且 \(\zeta_+\zeta_-=r^{-2}\)。

若 \(z\notin\mathcal E_r\)，则两根均不在单位圆上。
对充分大 \(z\) 有 \(\zeta_+\sim z/r\)、\(\zeta_-\sim 1/(rz)\)，故外部分支上恰有一根在单位圆内。
由连通性与根模长连续性，该性质在整个外部分支上保持。
于是外部情形中
\[
G_r(z)=\operatorname{Res}_{\zeta=\zeta_-}\frac{1}{D(\zeta)}
=\frac{1}{z-2r\zeta_-}
=\frac{1}{\sqrt{z^2-4}}.
\]

若 \(r\ne1\) 且 \(z\) 位于内部，则取 \(z=0\) 时 \(\zeta_\pm=\pm \mathrm{i}/r\) 均落在单位圆内。
同样由连通性可知，整个内部都保持两根在单位圆内。
因此内部留数和为
\[
\frac{1}{z-2r\zeta_-}+\frac{1}{z-2r\zeta_+}
=\frac{1}{\sqrt{z^2-4}}+\frac{1}{-\sqrt{z^2-4}}
=0,
\]
故 \(G_r(z)=0\)。证毕。
\end{proof}
\leanverified{paper\_prime\_register\_joukowsky\_cauchy\_rigidity}

\begin{corollary}[解析矩完全同一：中心二项系数律]\label{cor:prime-register-joukowsky-analytic-moments}
令解析矩
\[
m_n(r):=\int_{\mathcal E_r}w^n\,d\mu_r(w).
\]
则对一切 \(r>0\) 与 \(k\ge0\)，有
\[
m_{2k+1}(r)=0,
\qquad
m_{2k}(r)=\binom{2k}{k}.
\]
特别地，\(\{m_n(r)\}_{n\ge0}\) 与 \(r\) 无关。
\end{corollary}

\begin{proof}
对充分大 \(|z|\)，
\[
G_r(z)=\int\frac{1}{z-w}\,d\mu_r(w)
=\frac{1}{z}\sum_{n\ge0}\frac{m_n(r)}{z^n}.
\]
另一方面，由定理 \ref{thm:prime-register-joukowsky-cauchy-rigidity}，
\[
G_r(z)=\frac{1}{\sqrt{z^2-4}}
=\frac{1}{z}\frac{1}{\sqrt{1-4/z^2}}
=\frac{1}{z}\sum_{k\ge0}\binom{2k}{k}\frac{1}{z^{2k}}.
\]
比较系数即得结论。证毕。
\end{proof}
\leanverified{paper\_prime\_register\_joukowsky\_analytic\_moments}

\begin{proposition}[最小非解析见证：二阶径向矩反演尺度]\label{prop:prime-register-joukowsky-radial-moment-invert}
\leanverified{paper\_prime\_register\_joukowsky\_radial\_moment\_invert}
令
\[
R_2(r):=\int_{\mathcal E_r}|w|^2\,d\mu_r(w).
\]
则
\[
R_2(r)=r^2+r^{-2}.
\]
并且在规范 \(r\ge1\) 下，\(r\) 可由单一标量 \(R_2(r)\) 显式唯一恢复：
\[
\boxed{\;
r=\sqrt{\frac{R_2(r)+\sqrt{R_2(r)^2-4}}{2}}\; }.
\]
\end{proposition}

\begin{proof}
由 \(w=r e^{\mathrm{i}\theta}+r^{-1}e^{-\mathrm{i}\theta}\) 得
\[
|w|^2=r^2+r^{-2}+2\cos(2\theta).
\]
对 \(\theta\) 平均即得 \(R_2(r)=r^2+r^{-2}\)。
令 \(x=r^2\ge1\)，则 \(x+x^{-1}=R_2(r)\)。
解二次方程 \(x^2-R_2(r)x+1=0\) 即得所示闭式。证毕。
\end{proof}

\leanverified{paper\_prime\_register\_leyang\_joukowsky\_transport}
\begin{theorem}[Lee--Yang 单位圆零点的 Joukowsky--Gödel 椭圆输运]\label{thm:prime-register-leyang-joukowsky-transport}
设 \(K\subseteq\CC\) 为子域，\(P(z)\in K[z]\) 为次数 \(n\) 多项式，且其全部零点 \(\{z_j\}\subset\CC\) 满足 \(|z_j|=1\)。
对实参数 \(r>0\) 定义输运多项式
\[
Q_r(w):=\operatorname{Res}_{z}\!\bigl(P(z),\ rz^2-wz+r^{-1}\bigr)\in K(r)[w].
\]
则：
\begin{enumerate}
  \item \(Q_r(w)\) 的每个零点都落在 \(\mathcal E_r\) 上；
  \item 若 \(r\in K\)，则 \(Q_r(w)\in K[w]\)。
\end{enumerate}
\end{theorem}

\begin{proof}
由结式定义，\(Q_r(w)=0\) 当且仅当存在 \(z\in\CC\) 使
\[
P(z)=0,
\qquad
rz^2-wz+r^{-1}=0.
\]
第二式等价于 \(w=rz+r^{-1}z^{-1}\)。
若 \(|z|=1\)，则 \(z=e^{\mathrm{i}\theta}\)，从而 \(w=J_r(e^{\mathrm{i}\theta})\in\mathcal E_r\)，得 (1)。
(2) 则由结式是系数的多项式函数直接推出。证毕。
\end{proof}

\begin{corollary}[有理 Gödel 尺度的整系数闭包]\label{cor:prime-register-rational-godel-integer-closure}
\leanverified{paper\_prime\_register\_rational\_godel\_integer\_closure}
若 \(P(z)\in\ZZ[z]\)，且 \(r=a/b\in\QQ_{>0}\)（\(\gcd(a,b)=1\)），定义去分母输运多项式
\[
\widetilde Q_{a,b}(w)
:=  
\operatorname{Res}_{z}\!\bigl(P(z),\ a^2z^2-ab\,wz+b^2\bigr)\in\ZZ[w].
\]
则 \(\widetilde Q_{a,b}\) 的零点仍全部落在 \(\mathcal E_{a/b}\) 上，并且其构造仅依赖于 \(a\) 与 \(b\) 的普通素因子指数。
\end{corollary}

\begin{proof}
由 Sylvester 行列式表达式可知 \(\widetilde Q_{a,b}\in\ZZ[w]\)。
若 \(\widetilde Q_{a,b}(w)=0\)，则存在 \(|z|=1\) 使
\[
P(z)=0,
\qquad
a^2z^2-ab\,wz+b^2=0.
\]
于是
\[
w=\frac{a}{b}z+\frac{b}{a}z^{-1}=J_{a/b}(z)\in\mathcal E_{a/b}.
\]
最后一句来自唯一素因子分解。证毕。
\end{proof}

\begin{theorem}[椭圆化外置的可逆性：几何量确定素数尺度]\label{thm:prime-register-ellipse-invert-godel}
\leanverified{paper\_prime\_register\_ellipse\_invert\_godel}
在规范 \(r\ge1\) 下，下述任一类观测足以唯一恢复尺度 \(r\)，从而在 \(r\in\QQ_{>0}\) 时唯一恢复其素数指数向量 \(a\in\ZZ^{(\PP)}\)：
\begin{enumerate}
  \item 观测椭圆 \(\mathcal E_r\) 的半轴 \((A,B)\)，则 \(r=(A+B)/2\)；
  \item 观测推前测度 \(\mu_r\) 的二阶径向矩 \(R_2(r)\)，则 \(r\) 由命题 \ref{prop:prime-register-joukowsky-radial-moment-invert} 的闭式给出。
\end{enumerate}
\end{theorem}

\begin{proof}
(1) 由命题 \ref{prop:prime-register-godel-joukowsky-ellipse} 中 \(A(r)=r+r^{-1}\)、\(B(r)=r-r^{-1}\) 得 \(A+B=2r\)。
(2) 由命题 \ref{prop:prime-register-joukowsky-radial-moment-invert}。
最后，当 \(r\in\QQ_{>0}\) 时，由算术基本定理存在且仅存在有限支撑 \(a\in\ZZ^{(\PP)}\) 使 \(r=\mathsf G(a)\)。证毕。
\end{proof}

\paragraph{Gödel 对数尺度与端点 Busemann 时间}
\begin{proposition}[Gödel 对数尺度的 Busemann 时间同构]\label{prop:prime-register-godel-log-busemann-time}
沿用定义 \ref{def:prime-register-godel-scale-joukowsky} 的 Gödel 尺度 \(\mathsf{G}:\ZZ^{(\PP)}\to\QQ_{>0}\)，并定义对数时间参数
\[
t(a):=\log \mathsf{G}(a)\qquad(a\in\ZZ^{(\PP)}).
\]
则对任意 \(a,b\in\ZZ^{(\PP)}\) 有 \(t(a+b)=t(a)+t(b)\)。
并且对任意 \(u\in\mathbb{D}\setminus\{1\}\) 有
\[
y\!\bigl(\psi_{\mathsf{G}(a)}(u)\bigr)=y(u)+t(a),
\]
其中 \(\psi_m\) 与 \(y(u)\) 见 \S\ref{subsec:unit-circle-scale-linearization}（式 \eqref{eq:endpoint-busemann-addition-law}）。
更进一步，
\[
t(a)=\sum_{p\in\PP}a_p\log p.
\]
\leanverified{godelScale}
\leanverified{godelLogTime}
\leanverified{godelScale\_zero}
\leanverified{godelLogTime\_zero}
\leanverified{godelLogTime\_eq\_sum}
\leanverified{godelScale\_add}
\leanverified{godelLogTime\_add}
\leanverified{paper\_prime\_register\_godel\_log\_busemann\_time}
\end{proposition}
\begin{proof}
由 \(\mathsf{G}(a+b)=\mathsf{G}(a)\mathsf{G}(b)\) 得 \(t(a+b)=t(a)+t(b)\)。
将 \(m=\mathsf{G}(a)\) 代入 \eqref{eq:endpoint-busemann-addition-law} 即得第二式。
最后一式由 \(\mathsf{G}(a)=\prod_{p}p^{a_p}\) 取对数得到。证毕。
\end{proof}

\endinput



\begin{remark}[可观测正规形与外置账本的分解接口]\label{rem:prime-register-observable-vs-ledger}
定理 \ref{thm:prime-register-residual-ledger-group} 给出一条可直接核验的分解：\(\mathrm{NF}_{\mathrm{pr}}(r)\in\mathrm{ZG}\) 是由值不变量决定的可观测正规形，而 \(R(r)\in H\) 则携带被正规化投影压缩掉的区分信息。
二者合并为单射 \(\Theta(r)=(\mathrm{NF}_{\mathrm{pr}}(r),R(r))\)，从而把信息守恒落实为一份显式的因子分解与最小外置。
另一方面，定义 \ref{def:prime-register-godel-scale-joukowsky} 至定理 \ref{thm:prime-register-ellipse-invert-godel} 给出与该外置口径一致的复迹椭圆化接口：解析观测在外部支上对尺度 \(r\) 呈现刚性不可见性，而径向二阶矩给出单标量可逆读出。
\end{remark}

\subsubsection{Brauer 障碍的素数寄存器化：对角圆锥曲线、局部分歧与 XOR 合成}\label{subsubsec:prime-register-brauer-obstruction}
本段把一类可完全分类的中心单代数障碍转译为有限支撑的素数寄存器向量，并与“二次扩张足以消除障碍”的同余筛选判据对齐。

\begin{definition}[对角圆锥曲线与四元数 Brauer 类]\label{def:prime-register-brauer-diagonal-conic-quaternion}
取平方自由 \(a,b\in\QQ_{>0}\)，并考虑射影圆锥曲线
\[
C_{a,b}:\ aX^{2}+bY^{2}=Z^{2}\subset \mathbb{P}^{2}_{\QQ}.
\]
令 \(A_{a,b}:=(a,b)_{\QQ}\) 为相应的四元数代数（其 Brauer 类属于 \(\mathrm{Br}(\QQ)[2]\)）。
对每个素数 \(p\) 令
\[
\operatorname{inv}_p:\ \mathrm{Br}(\QQ_p)\longrightarrow \QQ/\ZZ
\]
为局部不变量，同样记
\(
\operatorname{inv}_\infty:\mathrm{Br}(\RR)\to \QQ/\ZZ
\)。
定义有限分歧集
\[
\operatorname{Ram}_{\mathrm{fin}}(A_{a,b})
:=\{\,p:\ \operatorname{inv}_p(A_{a,b})=\tfrac12\,\},
\qquad
\operatorname{Ram}(A_{a,b})
:=\operatorname{Ram}_{\mathrm{fin}}(A_{a,b})\cup\{\infty:\operatorname{inv}_\infty(A_{a,b})=\tfrac12\}.
\]
\end{definition}

\begin{theorem}[分歧素数集的有限性与偶性]\label{thm:prime-register-brauer-ramification-even}
对任意四元数代数 \(A/\QQ\)，集合 \(\operatorname{Ram}(A)\) 有限且基数为偶数。
若额外满足 \(\operatorname{inv}_\infty(A)=0\)（等价于 \(A\otimes_{\QQ}\RR\simeq \mathrm{Mat}_2(\RR)\)），则
\[
\bigl|\operatorname{Ram}_{\mathrm{fin}}(A)\bigr|\ \text{为偶数}.
\]
\leanverified{half\_sum\_eq\_card\_div\_two}
\leanverified{half\_sum\_int\_iff\_card\_even}
\leanverified{reciprocity\_forces\_even\_card}
\leanverified{even\_card\_gives\_int\_half\_sum}
\leanverified{paper\_prime\_register\_brauer\_ramification\_even}
\end{theorem}
\begin{proof}
四元数代数在每个地方 \(v\) 的局部不变量仅取 \(0\) 或 \(\tfrac12\)。
Brauer 群全局互反律给出
\(
\sum_{v}\operatorname{inv}_v(A)=0\in\QQ/\ZZ
\)，
其中求和遍历全部素数与无穷处，且除有限多个 \(v\) 外不变量为 \(0\)。
因此 \(\operatorname{inv}_v(A)=\tfrac12\) 的地方数必须为偶数。
在 \(\operatorname{inv}_\infty(A)=0\) 时，偶性落在有限素数部分。参见 \cite{Serre1973CourseArithmetic,VoightQuaternionAlgebras}。证毕。
\end{proof}

\begin{theorem}[Severi--Brauer 判据：有理点当且仅当分裂]\label{thm:prime-register-brauer-severi-brauer-splitting}
\leanverified{paper\_prime\_register\_brauer\_severi\_brauer\_splitting}
在定义 \ref{def:prime-register-brauer-diagonal-conic-quaternion} 的记号下，对任意域扩张 \(L/\QQ\)，有等价性
\[
C_{a,b}(L)\neq\varnothing
\quad\Longleftrightarrow\quad
A_{a,b}\otimes_{\QQ}L\simeq \mathrm{Mat}_2(L).
\]
\end{theorem}
\begin{proof}
圆锥曲线 \(C_{a,b}\) 是一个 Severi--Brauer 曲线，其 Brauer 类即为 \(A_{a,b}\) 的类。
Severi--Brauer 变体的基本判据断言：存在 \(L\)-点当且仅当对应中心单代数在 \(L\) 上分裂。
参见 \cite[Chapter~5]{GilleSzamuely2006CentralSimpleAlgebras} 或 \cite[Chapter~5]{Skorobogatov2001TorsorsRationalPoints}。证毕。
\end{proof}

\begin{proposition}[二次外置的局部分裂判据：分歧素数不允许分裂]\label{prop:prime-register-brauer-quadratic-splitting}
设 \(A/\QQ\) 为四元数代数，且 \(\operatorname{inv}_\infty(A)=0\)。
取平方自由 \(d\in\ZZ\setminus\{0\}\) 并令 \(L:=\QQ(\sqrt d)\)。
则
\[
A\otimes_{\QQ}L\simeq \mathrm{Mat}_2(L)
\quad\Longleftrightarrow\quad
\forall\,p\in \operatorname{Ram}_{\mathrm{fin}}(A),\ p\ \text{在}\ L\ \text{中不分裂}.
\]
\end{proposition}
\begin{proof}
对任意素数 \(p\)，有局部判据
\(
A\otimes_{\QQ}L\simeq \mathrm{Mat}_2(L)
\Longleftrightarrow
A\otimes_{\QQ}\QQ_p\ \text{在}\ L\otimes_{\QQ}\QQ_p\ \text{上分裂}
\)。
当 \(p\notin \operatorname{Ram}_{\mathrm{fin}}(A)\) 时，\(A\otimes\QQ_p\) 已分裂，故局部无约束。
当 \(p\in \operatorname{Ram}_{\mathrm{fin}}(A)\) 时，\(A\otimes\QQ_p\) 为除环，其最大交换子域恰为二次扩张；因此它在 \(\QQ_p(\sqrt d)\) 上分裂当且仅当 \(\QQ_p(\sqrt d)\) 为域，亦即 \(p\) 在 \(L\) 中不分裂。
参见 \cite[Chapter~12]{Serre1973CourseArithmetic} 或 \cite[Chapter~14]{VoightQuaternionAlgebras}。证毕。
\end{proof}
\leanverified{paper\_prime\_register\_brauer\_quadratic\_splitting}

\begin{corollary}[同余类筛选与 Chebotarev 密度]\label{cor:prime-register-brauer-congruence-sieve-density}
在命题 \ref{prop:prime-register-brauer-quadratic-splitting} 的设定下，令
\[
S:=\operatorname{Ram}_{\mathrm{fin}}(A),\qquad
M:=8\prod_{p\in S}p.
\]
则存在显式构造的剩余类集合 \(\mathcal{A}\subseteq(\ZZ/M\ZZ)^{\times}\) 使得对任意素数 \(\ell\nmid M\) 有
\[
A\otimes_{\QQ}\QQ(\sqrt{\ell})\simeq \mathrm{Mat}_2(\QQ(\sqrt{\ell}))
\quad\Longleftrightarrow\quad
\ell\bmod M\in \mathcal{A}.
\]
并且满足该条件的素数集合在全体素数中具有自然密度
\[
2^{-|S|}.
\]
\leanverified{single\_binary\_density}
\leanverified{two\_binary\_density}
\leanverified{three\_binary\_density}
\leanverified{independent\_binary\_density}
\leanverified{neg\_power\_eq\_reciprocal}
\leanverified{density\_multiplicative}
\leanverified{density\_lt\_one}
\leanverified{paper\_prime\_register\_brauer\_congruence\_sieve\_density}
\end{corollary}
\begin{proof}
对奇素数 \(p\in S\)，\(p\) 在 \(\QQ(\sqrt{\ell})\) 中不分裂等价于 Legendre 符号 \(\left(\frac{\ell}{p}\right)\neq 1\)，在 \(\ell\nmid p\) 情形下等价于 \(\left(\frac{\ell}{p}\right)=-1\)。
对 \(p=2\) 的条件可用 \(\ell\bmod 8\) 表述。
这些条件把 \(\ell\) 限制在模 \(M\) 的有限多个剩余类中，由中国剩余定理给出 \(\mathcal{A}\) 的存在与可构造性。
密度结论来自二次特征的 Chebotarev 等分布：每个 \(p\in S\) 对应一个独立的二次角色约束，使可行类数按 \(2^{-|S|}\) 比例收缩。参见 \cite[Chapter~6]{Skorobogatov2001TorsorsRationalPoints}。证毕。
\end{proof}

\begin{theorem}[有限分歧 Brauer 类的素数寄存器同构与 XOR 律]\label{thm:prime-register-brauer-xor-law}
令
\[
\mathcal{B}:=\{\, [A]\in \mathrm{Br}(\QQ)[2]:\ \operatorname{inv}_\infty(A)=0\,\}.
\]
定义映射
\[
\Psi:\ \mathcal{B}\longrightarrow \{\,S\subset \{\text{素数}\}:\ S\ \text{有限且}\ |S|\ \text{为偶数}\,\},
\qquad
[A]\longmapsto \operatorname{Ram}_{\mathrm{fin}}(A).
\]
则 \(\Psi\) 为交换群同构（右侧群运算为对称差 \(\Delta\)）。
特别地，对任意 \([A],[B]\in\mathcal{B}\)，有
\[
\operatorname{Ram}_{\mathrm{fin}}(A\otimes B)
=\operatorname{Ram}_{\mathrm{fin}}(A)\ \Delta\ \operatorname{Ram}_{\mathrm{fin}}(B).
\]
\leanverified{finset\_symmDiff\_comm}
\leanverified{finset\_symmDiff\_self}
\leanverified{finset\_symmDiff\_empty}
\leanverified{finset\_symmDiff\_assoc}
\leanverified{finset\_card\_symmDiff}
\leanverified{finset\_symmDiff\_card\_mod\_two}
\leanverified{finset\_symmDiff\_card\_even\_of\_even}
\leanverified{empty\_card\_even}
\leanverified{EvenCardFinset}
\leanverified{paper\_prime\_register\_brauer\_xor\_law}
\end{theorem}
\begin{proof}
对任意素数 \(p\)，局部不变量满足可加性
\(
\operatorname{inv}_p(A\otimes B)=\operatorname{inv}_p(A)+\operatorname{inv}_p(B)
\)。
由于四元数 \(2\)-挠类的取值仅为 \(0\) 或 \(\tfrac12\)，故
\(
\operatorname{inv}_p(A\otimes B)=\tfrac12
\)
当且仅当 \(\operatorname{inv}_p(A)\) 与 \(\operatorname{inv}_p(B)\) 恰有一个为 \(\tfrac12\)，这正是对称差规则。
互反律给出偶性约束（定理 \ref{thm:prime-register-brauer-ramification-even}）。
最后，四元数代数在 \(\QQ\) 上由其分歧集合完全决定，且任意有限偶素数集均可实现为某个 \(\operatorname{Ram}_{\mathrm{fin}}(A)\)（参见 \cite{VoightQuaternionAlgebras}），从而 \(\Psi\) 为同构。证毕。
\end{proof}

\begin{remark}[外置素数支持的最小性口径]\label{rem:prime-register-brauer-minimal-support}
在本段的圆锥曲线模型中，有限分歧集 \(S=\operatorname{Ram}_{\mathrm{fin}}(A_{a,b})\) 同时承担两类结构角色：
\begin{enumerate}
  \item 作为 Brauer 障碍的最小局部支持，它刻画了所有可能导致不可分裂的有限地方集合；
  \item 作为素数寄存器的有限支撑 \(\FF_2\)-向量，它把可合成算法类的群结构以 XOR 律外置到逐素数的可审计对象上（定理 \ref{thm:prime-register-brauer-xor-law}）。
\end{enumerate}
命题 \ref{prop:prime-register-brauer-quadratic-splitting} 与推论 \ref{cor:prime-register-brauer-congruence-sieve-density} 表明：在该可控子类中，二次扩张已足以把“分歧支持”转译为一个可算法化的同余筛选判据。
\end{remark}


